% Flavor catalog per-process template.
% process_id: E001
% family: edm_neutrino
% owner_agent_id: PKA-E001
% writer_agent_id: null
% checker_agent_id: null
% opus_signoff_id: null
% Canonical metadata sidecar: E001.yaml
% Status is controlled by the YAML sidecar status_history, not this comment.

\section{Electron electric dipole moment}

\subsection*{Process}
\textbf{Standard notation:} \(d_e\).  This entry covers the flavor-diagonal
CP-odd electron dipole moment, cataloged in the \texttt{edm\_neutrino} family
as an EDM-adjacent lepton-sector probe.  It is not a quark \(\Delta F=2\)
observable; its RS relevance comes from CP-violating lepton dipoles generated
by heavy lepton, neutrino, Higgs, or gauge-sector dynamics.

\subsection*{PDG / equivalent value(s)}
PDG Live 2026 lists the current bound as
\[
  |d_e| < 4.1\times 10^{-30}\ e\,\mathrm{cm}
  \qquad (90\%~\mathrm{CL}),
\]
from the Roussy et al. 2023 HfF\({}^+\) molecular-ion measurement; see
\texttt{flavor\_catalog/references/E001/pdg2026\_electron\_edm\_datablock.txt}.
The PDG datablock reports the table value as \(<0.041\) in units of
\(10^{-28}\ e\,\mathrm{cm}\), and records the corresponding measurement as
\((-1.3\pm2.0_{\rm stat}\pm0.6_{\rm syst})\times10^{-30}\ e\,\mathrm{cm}\).
The superseded ACME/Andreev 2018 ThO benchmark is
\(|d_e|<1.1\times10^{-29}\ e\,\mathrm{cm}\) at \(90\%\) CL, with the PDG
measurement note \((4.3\pm3.1\pm2.6)\times10^{-30}\ e\,\mathrm{cm}\).

\subsection*{Relevance to the RS / anarchic-flavor pipeline}
The electron EDM constrains the imaginary part of a CP-odd lepton dipole,
schematically \(\mathcal{L}\supset -i d_e\,\bar e\sigma_{\mu\nu}\gamma_5 e
F^{\mu\nu}/2\).  In warped lepton extensions, anarchic Yukawas and new CP
phases can feed this operator through loop matching.  The constraint is
therefore a direct diagnostic of CP alignment in the charged-lepton sector,
and it is especially relevant when the neutrino sector or bulk Higgs dynamics
introduce additional complex spurions.  It does not directly constrain the
current quark-only \(\Delta F=2\) scan.

\subsection*{Post-2008 developments}
Relative to the 2008 CFW RS-flavor reference point, the experimental EDM
frontier changed qualitatively.  ACME 2018 set the previous molecular ThO
limit at \(1.1\times10^{-29}\ e\,\mathrm{cm}\), already strong enough to
motivate EFT treatments of CP-violating electron dipoles.  Roussy et al. 2023
then improved the previous best upper bound by about \(2.4\) using trapped
HfF\({}^+\) molecular ions, yielding the current PDG value above.  The
post-ACME EFT literature, represented here by Panico, Pomarol, and Riembau
2019, treats electron EDMs as constraints on CP-violating dipole operators
and their two-loop anomalous dimensions, which is the natural language for a
future Wilson-coefficient implementation.

\subsection*{Constraint validity and model dependence}
The experimental null limit is robust as a bound on the effective electron EDM
extracted from paramagnetic molecular spectroscopy.  Model interpretation is
less automatic: paramagnetic systems can also be sensitive to semileptonic
CP-odd electron-nucleon interactions, and an RS implementation must specify
which SMEFT or low-energy CP-odd operator basis is being matched.  As a
catalog constraint this is a lepton-extension and EDM-adjacent loop bound, not
a standalone exclusion in the existing quark pipeline.

\subsection*{Code coverage in this repo}
\textbf{NO.}  The required catalog greps and a focused implementation grep for
\texttt{electron.?EDM}, \texttt{eEDM}, \texttt{electric dipole}, and
\texttt{d\_e} over \texttt{quarkConstraints/}, \texttt{qcd/},
\texttt{flavorConstraints/}, \texttt{neutrinos/}, \texttt{yukawa/},
\texttt{warpConfig/}, \texttt{solvers/}, \texttt{scanParams/}, and
\texttt{tests/} found no electron-EDM implementation.  The only adjacent
dipole code is the off-diagonal \(\mu\to e\gamma\) bound in
\texttt{flavorConstraints/muToEGamma.py:1}, \texttt{:3},
\texttt{:20--21}, and \texttt{:81}; \texttt{scanParams/scan.py:32--33}
sets a MEG II branching-fraction limit for that different process.

\subsection*{Implementation difficulty}
\textbf{HIGH.}  A live implementation needs a new flavor-diagonal CP-odd
lepton dipole observable, matching from the model's lepton/KK sector to the
dipole Wilson coefficient, and likely SMEFT-to-low-energy running.  The
existing \(\Delta F=2\) operator basis and the \(\mu\to e\gamma\) NDA helper
do not provide the needed CP-odd \(d_e\) calculation.

\subsection*{Key references}
Process-local reference keys before bibliography consolidation:
\texttt{PDG2026:ElectronEDM}, \texttt{Roussy:ElectronEDM2023},
\texttt{ACME:ElectronEDM2018},
\texttt{PanicoPomarolRiembau:ElectronEDM2019}, and
\texttt{CsakiFalkowskiWeiler:RSFlavor2008}.
